\chapter*{Preface to Part II}
\markboth{\MakeUppercase{Preface to Part II}}{\MakeUppercase{Preface to Part II}}
\addcontentsline{toc}{chapter}{Preface to Part II}
Since we have already given a general outline of the aim and scope of these notes in the preface to \emph{Several Complex Variables and Complex Manifolds I}, we shall do no more here than provide a brief description of the contents of this volume together with a few notes of guidance to the reader.

Chapter 5 of the present notes is devoted to calculus on complex manifolds. The first four sections cover basic linear algebra and calculus on a differential manifold. Most of the material in these sections should be familiar (though perhaps not the notation) and I would suggest reading through them quickly, referring back to them later, if necessary, for specific results and notation. The next three sections lead up to the construction of the $\bar\partial$-operator on an arbitrary complex manifold and also describe the $\bar\partial$-operator for holomorphic bundle valued forms. In section 8 we prove the Dolbeault--Grothendieck lemma and solve the Cousin problems for polydiscs in $\mathbb{C}^n$. In section 9 we show how holomorphic vector bundles naturally enter into the study of compact complex manifolds. We discuss, for example, the Euler sequence for projective space; the geometric genus; theta functions for complex tori. Finally in section 10, we discuss various pseudoconvexivity conditions for non-compact complex manifolds.

Chapter 6 is a self-contained introduction to the theory of sheaves in complex analysis. Section 1 is devoted to sheaves and presheaves with many examples: In section 2 we show how sheaf theory can be used to construct the envelope of holomorphy of a domain spread in $\mathbb{C}^n$. This section is not used elsewhere in the text and may be omitted at first reading. In section 3 we define sheaf cohomology using fine resolutions. After proving Leray's theorem, we go on to define \v{C}ech cohomology and prove that it is naturally isomorphic to cohomology computed using fine resolutions. There are many important illustrations of sheaf cohomology arguments in this section. For example: The de Rham isomorphism between singular and de Rham cohomology; The Dolbeault isomorphism theorem; the first Chern class and classification of complex line bundles.

In Chapter 7 we prove a number of foundational results in the theory of complex manifolds. In section 1 we define coherence and prove Oka's theorem on the coherence of the sheaf of relations. In
% Source PDF page 7; printed page vi.
section 2 we prove Cartan's theorems A and B granted the exactness of the $\bar\partial$-sequence for locally free sheaves (A proof of this result will be included in the projected part III of these notes; proofs may also be found in H\"ormander \cite{bib:II:hormander-l:1} or Vesentini \cite{bib:II:vesentini-e:1}). In section 3 we prove the finiteness theorem of Cartan and Serre and in section 4 the finiteness theorem of Grauert. In section 5 we prove Serre's theorems A and B and give a number of applications. In section 6 we prove Grauert's vanishing theorem and, following Grauert, show how it may be used to prove the Kodaira embedding theorem.

